% Part: second-order-logic
% Chapter: syntax-and-semantics
% Section: semantic-notions

\documentclass[../../../include/open-logic-section]{subfiles}

\begin{document}

\olfileid{sol}{syn}{sem}
\olsection{Gagasan Semantis}

\begin{explain}
Gagasan logis utama babagan \emph{validitas}, \emph{akibat
semantis}, lan \emph{bisa dicukupi} ditetepake kanthi cara kang
padha ing logika orde kapindho lan logika orde kapisan, kejaba
relasi nyukupi kang dadi dhasare saiki kanggo !!{formula}s
orde kapindho. Mesthi wae, !!{sentence} orde kapindho iku
!!a{formula} kang kabeh variabele, kalebu variabel predikat
lan fungsi, kaiket.
\end{explain}

\begin{defn}[Validitas]
Ukara $!A$ iku \emph{valid}, $\Entails !A$, yen lan mung yen
$\Sat{M}{!A}$ kanggo saben !!{structure}~$\Struct M$.
\end{defn}

\begin{defn}[Akibat Semantis]
Himpunan ukara~$\Gamma$ \emph{ngakibatake} ukara~$!A$,
$\Gamma \Entails !A$, yen lan mung yen kanggo saben
!!{structure}~$\Struct M$ kang ndadekake
$\Sat{M}{\Gamma}$, uga $\Sat{M}{!A}$.
\end{defn}

\begin{defn}[Bisa Dicukupi]
Himpunan ukara~$\Gamma$ \emph{bisa dicukupi} yen
$\Sat{M}{\Gamma}$ kanggo sawenehing
!!{structure}~$\Struct M$. Yen $\Gamma$ ora bisa dicukupi,
himpunan kasebut diarani \emph{ora bisa dicukupi}.
\end{defn}

\end{document}
